% !TEX root = ../../hw02.tex
% Homework 02, Exercise 2. Shared solution.

\begin{proof}[Solution]
Let
\[
  I_n=[n,\infty),\qquad n\geq1.
\]
Each $I_n$ is a nonempty closed interval in $\R$, since $n\in I_n$
and its complement $(-\infty,n)$ is open. Moreover,
\[
  I_{n+1}=[n+1,\infty)\subset[n,\infty)=I_n,
\]
so the intervals are nested; see
\autoref{fig:hw02:unbounded-intervals}.

For any $x\in\R$, the Archimedean property gives a positive
integer $n>x$.
Then $x\notin[n,\infty)=I_n$. Consequently, no real number belongs
to every $I_n$, and
\[
  \bigcap_{n=1}^{\infty}I_n=\varnothing.
\]

\begin{marginfigure}
  \centering
  % !TEX root = ../../problems/hw02.tex
% [n,infinity), n = 1,2,3,...; infinity is a direction, not an endpoint.
\begin{tikzpicture}[interval diagram]
  \foreach \n in {1,2,3} {
    \coordinate (left) at ({0.7*(\n-1)},{1-\n});
    \draw[->] (left) -- (3.6,{1-\n});
    \node[interval closed] at (left) {};
    \node[interval label,above=3pt] at (left) {$\n$};
  }
  \node[interval label,above=3pt] at (3.6,0) {$+\infty$};
  \node at (1.8,-2.9) {$\vdots$};
\end{tikzpicture}

  \caption[Nested unbounded intervals.]{Each ray $[n,\infty)$ includes its finite left endpoint
    and extends indefinitely to the right.}
  \label{fig:hw02:unbounded-intervals}
\end{marginfigure}

The missing hypothesis in the nested interval property is
\emph{boundedness}. All the intervals $I_n$ are unbounded, and
their left endpoints eventually exceed every fixed real number.
\end{proof}
